\subsubsection{Bellmanford}

Grafo direcionado com pesos, $n$ vértices, $m$ arestas.
Para um vértice especificado, diz a menor distância de todos os
outros até ele. Funciona com peso negativo.

Claramente, se tiver ciclo negativo não existe esse caminho. Então
dá pra detectar ciclo negativo também.

Roda em $O(nm)$

\inccode{grafos/bellmanford.cpp}

\subsubsection{SPFA}

Melhoramento do Bellmanford que toma vantagem do fato de que nem todas
as tentativas de relaxação vão funcionar. Também dá pra usar pra detectar
ciclo negativo.

O pior caso é $O(nm)$ também. Acontece que, na prática, é bem mais rápido, apesar
de ser fácil de criar contra-exemplos em que roda em $O(nm)$. Mas tem gente que diz
que em média roda em $O(m)$.

\inccode{grafos/spfa.cpp}